\section{Scalar Wave Equation In $3+1$ Form}

As a simple example of a $3+1$ decomposition, we first consider the scalar wave equation. This example illustrates the basic idea of rewriting a covariant equation as an initial value problem suitable for numerical time evolution.

We begin with the scalar wave equation with a source term,
\begin{equation}
\square \psi = \nabla_a \nabla^a \psi = 4\pi \rho .
\end{equation}
This equation is written in covariant form. In flat spacetime, using the Minkowski metric with signature $(-,+,+,+)$, it becomes
\begin{equation}
\square \psi
=
\eta^{ab}\partial_a\partial_b \psi
=
\left(-\partial_t^2 + D^2\right)\psi ,
\end{equation}
where $D^2$ denotes the spatial Laplacian. Thus,
\begin{equation}
\partial_t^2 \psi = D^2 \psi - 4\pi \rho .
\end{equation}

To write this as a first order system in time, we introduce
\begin{equation}
\Pi = \partial_t \psi .
\end{equation}
The wave equation can then be written as
\begin{align}
\partial_t \psi &= \Pi , \\
\partial_t \Pi &= D^2 \psi - 4\pi \rho .
\end{align}
This form makes the time evolution structure explicit: given initial data for $\psi$ and $\Pi$, the system determines their future evolution.

This scalar field system provides a useful contrast with the systems considered later. The scalar wave equation has a well posed initial value problem and contains no constraint equations. In particular, the initial data for $\psi$ and $\Pi$ can be freely specified, subject only to suitable smoothness and boundary conditions. This is not the case for Maxwell's equations or general relativity, where the initial data must also satisfy constraint equations.

\section{Electricity and Magnetism}

The scalar wave equation provides a simple example of an evolution system with no constraints on the initial data. We now consider Maxwell's equations, which provide the next step in complexity: the system naturally separates into evolution equations and constraint equations. This example illustrates, in a familiar setting, the same basic structure that appears later in general relativity.

We present the decomposition in flat spacetime, following an approach similar to that of \textit{Numerical Relativity: Starting from Scratch} \cite{BaumgarteShapiro2021}. Working in flat spacetime allows us to emphasize the basic evolution and constraint structure before introducing the additional geometric terms that arise in curved spacetime.

We begin by introducing the Faraday tensor in terms of the electromagnetic vector potential $A_a$,
\begin{equation}
F_{ab} = \nabla_a A_b - \nabla_b A_a .
\end{equation}
This definition makes the antisymmetry of the Faraday tensor explicit,
\begin{equation}
F_{ab} = -F_{ba}.
\end{equation}
It also ensures that the homogeneous Maxwell equation,
\begin{equation}
\nabla_{[a}F_{bc]}=0 ,
\end{equation}
is satisfied automatically.

The remaining Maxwell equation is
\begin{equation}
\nabla_b F^{ab} = 4\pi j^a ,
\label{eq:maxwell_equation}
\end{equation}
where $j^a$ is the four current. In flat spacetime, we write its components schematically as
\begin{equation}
j^a = (\rho,J^i),
\end{equation}
where $\rho$ is the charge density and $J^i$ is the spatial current. Charge conservation is expressed by the continuity equation
\begin{equation}
\nabla_a j^a = 0 .
\end{equation}

We now decompose the vector potential into an electromagnetic scalar potential $\Phi$ and a spatial vector potential $A_i$. With metric signature $(-,+,+,+)$, we write this schematically as
\begin{equation}
A_a = (-\Phi,A_i).
\end{equation}
Using $E_i=F_{i0}$, we have
\begin{equation}
E_i
=
F_{i0}
=
\nabla_i A_0 - \nabla_0 A_i .
\end{equation}
In flat spacetime, this becomes
\begin{equation}
E_i
=
-\partial_t A_i - D_i\Phi ,
\label{eq:electric_field_potential}
\end{equation}
where $D_i$ denotes the spatial derivative. Taking the divergence gives
\begin{equation}
D_i E^i = -\partial_t D_i A^i - D_iD^i\Phi .
\end{equation}
The time component of Eq.~\eqref{eq:maxwell_equation} gives
\begin{equation}
D_i E^i = 4\pi \rho ,
\label{eq:gauss_law}
\end{equation}
which is Gauss's law. This equation is a constraint equation: it restricts the allowed initial data for the electric field and charge density, rather than providing an independent time evolution equation.

Equation~\eqref{eq:electric_field_potential} can also be rearranged to give
\begin{equation}
\partial_t A_i = -E_i - D_i\Phi .
\label{eq:vector_potential_evolution}
\end{equation}
This equation gives the time evolution of the spatial vector potential once the scalar potential, or equivalently a gauge condition, has been chosen.

We now consider the spatial components of Maxwell's equations. From the definition
\begin{equation}
B^i = \frac{1}{2} \, \epsilon^{ijk}F_{jk},
\end{equation}
and the relation $F_{ij}=D_iA_j-D_jA_i$, the magnetic field can be written in terms of the vector potential as
\begin{equation}
B^i = \epsilon^{ijk}D_j A_k .
\end{equation}
This form automatically satisfies the magnetic constraint in flat spacetime,
\begin{equation}
D_i B^i = 0 .
\end{equation}

The spatial components of Eq.~\eqref{eq:maxwell_equation} give the evolution equation for the electric field. In tensor component form, this is Ampere's law,
\begin{equation}
\partial_t E_i
=
\epsilon_i{}^{jk}D_j B_k
-4\pi J_i ,
\label{eq:ampere_tensor}
\end{equation}
where $\epsilon_i{}^{jk}$ is the spatial Levi Civita tensor. Substituting $B^i=\epsilon^{ijk}D_jA_k$ gives the equivalent potential form
\begin{equation}
\partial_t E_i
=
-D_jD^j A_i
+ D_iD^j A_j
- 4\pi J_i .
\label{eq:electric_field_evolution_potential}
\end{equation}

Thus, in this flat spacetime example, Maxwell's equations separate into evolution equations,
\begin{align}
\partial_t A_i &= -E_i - D_i\Phi , \\
\partial_t E_i
&=
-D_jD^j A_i
+ D_iD^j A_j
- 4\pi J_i ,
\end{align}
and constraint equations,
\begin{align}
D_i E^i &= 4\pi \rho , \\
D_i B^i &= 0 .
\end{align}
The constraints restrict the allowed initial data, while the evolution equations determine how that data changes in time. If the constraints are satisfied initially and the continuity equation holds, the Maxwell evolution equations preserve the constraints.

This illustrates the same basic structure that appears in the $3+1$ decomposition of general relativity. Part of the system provides evolution equations, while the remaining equations constrain the allowed initial data. In general relativity, the analogous decomposition involves the spacetime metric, lapse, shift, spatial metric, extrinsic curvature, and spatial covariant derivative. The conservation laws are also tied to the geometric structure of the theory through the Bianchi identities.

\section{Spacetime and Kinematics}
Before proceeding further, we briefly discuss the metric in more detail. The metric encodes the notion of distance and separation in spacetime. As a simple example, the Minkowski metric describes flat spacetime. In Cartesian coordinates, and using the signature $(-,+,+,+)$, the line element is
\begin{equation}
ds^2 = -dt^2 + dx^2 + dy^2 + dz^2 .
\end{equation}
Equivalently, the metric components can be written in matrix form as
\begin{equation}
\eta_{ab} =
\begin{pmatrix}
-1 & 0 & 0 & 0 \\
0 & 1 & 0 & 0 \\
0 & 0 & 1 & 0 \\
0 & 0 & 0 & 1
\end{pmatrix},
\end{equation}
where the rows and columns are ordered as $(t,x,y,z)$. The negative sign in the time component reflects our choice of signature.

The same flat spacetime can be expressed in different coordinate systems. For example, in spherical coordinates the Minkowski line element becomes
\begin{equation}
ds^2 = -dt^2 + dr^2 + r^2 d\theta^2 + r^2\sin^2\theta\, d\phi^2 .
\end{equation}
In these coordinates, the metric components are
\begin{equation}
\eta_{ab} =
\begin{pmatrix}
-1 & 0 & 0 & 0 \\
0 & 1 & 0 & 0 \\
0 & 0 & r^2 & 0 \\
0 & 0 & 0 & r^2 \sin^2\theta
\end{pmatrix},
\end{equation}
where the rows and columns are ordered as $(t,r,\theta,\phi)$. Although these representations look different, they describe the same flat spacetime. This illustrates that the components of the metric depend on the choice of coordinates.

Before discussing curvature, we recall the definition of the covariant derivative and the Christoffel symbols. For scalars, the ordinary partial derivative is sufficient. For tensors, however, the derivative must also account for changes in the basis vectors. This is done by introducing the connection coefficients, or Christoffel symbols.

For a vector $A^b$, the covariant derivative is
\begin{equation}
\nabla_a A^b
=
\partial_a A^b
+
\Gamma^b{}_{ac}A^c .
\end{equation}
For a covector $A_b$, the corresponding expression is
\begin{equation}
\nabla_a A_b
=
\partial_a A_b
-
\Gamma^c{}_{ab}A_c .
\end{equation}
For the Levi-Civita connection, which is metric compatible and torsion free, the Christoffel symbols are
\begin{equation}
\Gamma^a{}_{bc}
=
\frac{1}{2}g^{ad}
\left(
\partial_b g_{dc}
+
\partial_c g_{db}
-
\partial_d g_{bc}
\right).
\end{equation}

A natural question is how to determine whether a given metric describes flat spacetime. A necessary and sufficient local condition is that the Riemann curvature tensor vanishes everywhere. Thus, curvature provides a way, independent of coordinates, to distinguish flat spacetime from genuinely curved spacetime.

 We will now build up the mathematical foundation we need to construct the so-called ADM decomposition and resulting ADM equations.  (ADM stands for Arnowitt--Deser--Misner, after Richard Arnowitt, Stanley Deser, and Charles Misner. The ADM equations emerge from just such a 3+1 decomposition of the Einstein equations.) This freedom {\bf EWH:  What freedom are you referring to? -- and do you wnat to reduce the previous parenthetical to a footnote?  Would that be better or worse?} is needed to formulate the Cauchy problem in general relativity and in alternative theories of gravity. We foliate the spacetime into a family of spacelike hypersurfaces $\Sigma_t$, which may be interpreted as slices of constant coordinate time $t$.

The gradient of the time function is
\begin{equation}
\nabla_a t = \partial_a t .
\end{equation}
This covector is normal to the constant time hypersurfaces. We define the future directed unit normal $n_a$ to the slices by
\begin{equation}
n_a = -\alpha \nabla_a t ,
\end{equation}
where $\alpha$ is the lapse function. The normalization condition is
\begin{equation}
n_a n^a = g^{ab}n_a n_b = -1 .
\end{equation}
Because the hypersurfaces are spacelike, their normal is timelike, and we therefore normalize it to negative unity. The lapse $\alpha$ measures the rate at which proper time advances relative to coordinate time for observers moving normal to the slices.

In coordinates adapted to the foliation,
\begin{equation}
\nabla_a t = \partial_a t = \frac{\partial t}{\partial x^a},
\end{equation}
so the only nonzero component is the time component. Therefore,
\begin{equation}
n_a = (-\alpha,0,0,0).
\end{equation}
For an observer moving normal to the slice, the proper time interval satisfies
\begin{equation}
d\tau = \alpha\,dt .
\end{equation}

The coordinate time evolution vector $t^a$ can be decomposed into a part normal to the slice and a part tangent to the slice:
\begin{equation}
t^a = \alpha n^a + \beta^a .
\end{equation}
Here $\beta^a$ is the shift vector. It is purely spatial, so
\begin{equation}
n_a\beta^a=0 .
\end{equation}
Equivalently,
\begin{equation}
\alpha n^a = t^a - \beta^a .
\label{eq:normal}
\end{equation}
In coordinates adapted to the foliation, $\beta^a=(0,\beta^i)$ and $t^a=(1,0,0,0)$. Therefore,
\begin{equation}
n^a = \frac{1}{\alpha}(1,-\beta^i).
\end{equation}
The lapse controls the separation between neighboring slices in proper time, while the shift describes how the spatial coordinates move from one slice to the next.

With these definitions, the spacetime line element can be written in ADM form as
\begin{equation}
ds^2
=
-\alpha^2 dt^2
+
\gamma_{ij}
\left(dx^i+\beta^i dt\right)
\left(dx^j+\beta^j dt\right).
\end{equation}
This form makes explicit the split between the time direction, described by $\alpha$ and $\beta^i$, and the spatial geometry, described by the spatial metric $\gamma_{ij}$.

The spacetime metric can be decomposed into spatial and temporal parts using the normal vector. The spatial metric is
\begin{equation}
\gamma_{ab} = g_{ab} + n_a n_b ,
\end{equation}
and the corresponding projection operator is
\begin{equation}
\gamma^a{}_{b} = \delta^a{}_{b} + n^a n_b .
\end{equation}
This operator projects spacetime tensors onto the spatial hypersurfaces.

The way the spatial hypersurfaces are embedded in spacetime is described by the extrinsic curvature. With the sign convention used in this work, we define
\begin{equation}
K_{ab}
=
-\gamma_a{}^c\gamma_b{}^d\nabla_c n_d .
\end{equation}
Equivalently, $K_{ab}$ measures the spatially projected change of the normal vector from point to point on the hypersurface.

The Lie derivative measures how a tensor field changes as it is dragged along the flow generated by a vector field. For a covector $A_d$, the Lie derivative along $n^a$ is
\begin{equation}
\mathcal{L}_n A_d
=
n^c\nabla_c A_d
+
A_c\nabla_d n^c .
\end{equation}
In the $3+1$ decomposition, the relation
\begin{equation}
t^a = \alpha n^a + \beta^a
\end{equation}
allows coordinate time evolution to be separated into evolution normal to the slice and spatial transport along the shift.

\section{General Relativity}

We now apply the $3+1$ decomposition introduced above to Einstein's equations. The fundamental dynamical field is the spacetime metric $g_{ab}$, which satisfies
\begin{equation}
G_{ab} + \Lambda g_{ab} = 8\pi T_{ab}.
\label{eq:Einstein}
\end{equation}
For the remainder of this discussion we set $\Lambda=0$, as is appropriate for the asymptotically flat spacetimes considered in this work.

At first, Einstein's equation appears to provide ten equations for the ten independent components of the spacetime metric. However, these equations are not all independent evolution equations. The Einstein tensor satisfies the contracted Bianchi identity,
\begin{equation}
\nabla_a G^{ab} = 0 .
\end{equation}
Together with Einstein's equation, this implies the local conservation law
\begin{equation}
\nabla_a T^{ab} = 0 .
\end{equation}
In the $3+1$ decomposition, this structure appears as four constraint equations together with evolution equations for the spatial metric and extrinsic curvature. This is analogous to the way Maxwell's equations contain both evolution equations and constraints, but in general relativity the constraints restrict the geometry of the initial spatial slice.

Using the foliation defined above, the spacetime metric is decomposed into spatial and temporal parts according to
\begin{equation}
\gamma_{ab} = g_{ab} + n_a n_b ,
\end{equation}
where $n^a$ is the future directed unit normal to the spatial hypersurfaces. Equivalently,
\begin{equation}
g_{ab} = \gamma_{ab} - n_a n_b .
\end{equation}
The tensor $\gamma_{ab}$ acts as the intrinsic metric on each spatial slice. Since it is purely spatial, its contraction with the normal vanishes:
\begin{equation}
\gamma_{ab}n^b = 0 .
\end{equation}
The inverse spacetime metric can similarly be written as
\begin{equation}
g^{ab} = \gamma^{ab} - n^a n^b .
\end{equation}

The extrinsic curvature relates the embedding of the spatial slice to the time evolution of the spatial metric. With the sign convention adopted in this work,
\begin{equation}
K_{ab}
=
-\frac{1}{2}\mathcal{L}_n \gamma_{ab}.
\label{eq:extrinsic-curvature-lie}
\end{equation}
Using the decomposition $t^a=\alpha n^a+\beta^a$, this gives
\begin{equation}
\partial_t \gamma_{ij}
=
-2\alpha K_{ij}
+
\mathcal{L}_{\beta}\gamma_{ij}.
\label{eq:adm-gamma-evolution}
\end{equation}
Thus, the extrinsic curvature measures how the spatial metric changes in time, up to the lapse and shift. The spatial metric $\gamma_{ij}$ and the extrinsic curvature $K_{ij}$ form the basic dynamical variables in the ADM decomposition, while $\alpha$ and $\beta^i$ encode coordinate freedom.

We also introduce the spatial covariant derivative $D_i$, which is compatible with the spatial metric $\gamma_{ij}$. For a spatial vector $V^j$,
\begin{equation}
D_i V^j
=
\partial_i V^j
+
{}^{(3)}\Gamma^j{}_{ik} V^k ,
\end{equation}
where
\begin{equation}
{}^{(3)}\Gamma^i{}_{jk}
=
\frac{1}{2}\gamma^{il}
\left(
\partial_j \gamma_{lk}
+
\partial_k \gamma_{lj}
-
\partial_l \gamma_{jk}
\right)
\end{equation}
is the Christoffel symbol associated with $\gamma_{ij}$.

The projections of the spacetime curvature are related to the intrinsic and extrinsic geometry of the spatial slices through the Gauss--Codazzi relations. These relations provide the geometric bridge between the four dimensional spacetime curvature and the three dimensional geometry of each spatial hypersurface.

The Gauss equation is
\begin{equation}
\gamma^a{}_{i}
\gamma^b{}_{j}
\gamma^c{}_{k}
\gamma^d{}_{l}
{}^{(4)}R_{abcd}
=
{}^{(3)}R_{ijkl}
+
K_{ik}K_{jl}
-
K_{il}K_{jk}.
\end{equation}
Here ${}^{(4)}R_{abcd}$ is the four dimensional spacetime Riemann tensor, while ${}^{(3)}R_{ijkl}$ is the three dimensional Riemann tensor associated with $\gamma_{ij}$.

The Codazzi equation is
\begin{equation}
\gamma^a{}_{i}
\gamma^b{}_{j}
\gamma^c{}_{k}
n^d
{}^{(4)}R_{abcd}
=
D_j K_{ik}
-
D_i K_{jk}.
\end{equation}
This relates the mixed spatial normal projection of the spacetime curvature to spatial derivatives of the extrinsic curvature.

The Ricci equation, which involves two projections along the normal direction, is
\begin{equation}
\gamma^a{}_{i}
\gamma^b{}_{j}
n^c n^d
{}^{(4)}R_{acbd}
=
\mathcal{L}_n K_{ij}
+
\frac{1}{\alpha}D_iD_j\alpha
+
K_{ik}K^k{}_{j}.
\end{equation}

To obtain the ADM equations, we insert these decompositions into Einstein's equation and equate the projections of the Einstein tensor with the corresponding projections of the stress energy tensor in Eq.~\eqref{eq:Einstein}.

Projecting Einstein's equation twice along the normal direction gives the Hamiltonian constraint:
\begin{equation}
R + K^2 - K_{ij}K^{ij}
=
16\pi \rho .
\label{eq:hamiltonian}
\end{equation}
Here $R$ is the scalar curvature associated with $\gamma_{ij}$,
\begin{equation}
K = \gamma^{ij}K_{ij}
\end{equation}
is the trace of the extrinsic curvature, and
\begin{equation}
\rho = n^a n^b T_{ab}
\end{equation}
is the energy density measured by an observer moving normal to the spatial slice.

A mixed projection, with one index projected along the normal direction and one projected onto the spatial slice, gives the momentum constraint:
\begin{equation}
D_j \left( K^{ij} - \gamma^{ij}K \right)
=
8\pi j^i .
\label{eq:momentum}
\end{equation}
Here
\begin{equation}
j^i
=
-\gamma^{ia}n^b T_{ab}
\end{equation}
is the momentum density measured by the normal observer.

A complete spatial projection gives the evolution equation for the extrinsic curvature:
\begin{equation}
\partial_t K_{ij}
=
\alpha
\left(
R_{ij}
- 2K_{ik}K^k{}_{j}
+ K K_{ij}
\right)
- D_iD_j\alpha
- 4\pi\alpha
\left[
2S_{ij}
-
\gamma_{ij}\left(S-\rho\right)
\right]
+
\mathcal{L}_{\beta}K_{ij}.
\label{eq:adm-k-evolution}
\end{equation}
Here $R_{ij}$ is the Ricci tensor associated with $\gamma_{ij}$, 
$S_{ij} = \gamma^a{}_{i}\gamma^b{}_{j}T_{ab}$ 
is the spatial projection of the stress energy tensor, and
$S = \gamma^{ij}S_{ij}$ is its trace.

Together with Eq.~\eqref{eq:adm-gamma-evolution}, these equations form the ADM system. The variables $\gamma_{ij}$ and $K_{ij}$ are the dynamical variables, while $\alpha$ and $\beta^i$ encode the gauge freedom associated with the choice of spacetime coordinates.

\begin{equation}
\boxed{
\begin{alignedat}{3}
R + K^2 - K_{ij}K^{ij}
&= 16\pi \rho,
&\qquad&
\text{Hamiltonian constraint},
\\[0.5em]
D_j \left( K^{ij} - \gamma^{ij}K \right)
&= 8\pi j^i,
&&
\text{Momentum constraint},
\\[0.5em]
\partial_t \gamma_{ij}
&= -2\alpha K_{ij}
+ \mathcal{L}_{\beta}\gamma_{ij},
&&
\text{Metric evolution},
\\[0.5em]
\partial_t K_{ij}
&= \alpha
\left(
R_{ij}
- 2K_{ik}K^k{}_{j}
+ K K_{ij}
\right)
- D_iD_j\alpha
&&
\\
&\quad
- 4\pi\alpha
\left[
2S_{ij}
-
\gamma_{ij}(S-\rho)
\right]
+ \mathcal{L}_{\beta}K_{ij},
&&
\text{Extrinsic curvature evolution}.
\end{alignedat}
}
\end{equation}
\subsection{Reformulating the Einstein equations}
The ADM equations provide a natural $3+1$ decomposition of Einstein's equations, but the standard ADM system is only weakly hyperbolic and is not generally well suited for stable numerical evolution. It is therefore useful to reformulate the equations into a system with better hyperbolic properties. The Baumgarte--Shapiro--Shibata--Nakamura (BSSN) formulation achieves this by introducing conformal variables and promoting certain combinations of spatial derivatives of the metric to independent evolved quantities \cite{Shibata1995_BSSN,Baumgarte1999_BSSN,Baumgarte2010}.

\subsubsection{A Maxwell analogy}

We first illustrate the basic idea with an analogy from electromagnetism. In a flat spacetime $3+1$ decomposition, Maxwell's equations can be written in terms of the electric field $E_i$ and vector potential $A_i$ as
\begin{equation}
\partial_t E_i
=
-D_jD^j A_i
+
D_iD^j A_j
-
4\pi J_i ,
\end{equation}
together with
\begin{equation}
\partial_t A_i
=
-E_i
-
D_i\Phi .
\end{equation}
Here $\Phi$ is the electromagnetic scalar potential, and $J_i$ is the spatial current. These evolution equations are supplemented by Gauss's law,
\begin{equation}
D_iE^i=4\pi\rho .
\end{equation}
We define the electric constraint violation by
\begin{equation}
\mathcal{C}
\equiv
D_iE^i-4\pi\rho .
\end{equation}
For constraint satisfying data, $\mathcal{C}=0$. In exact arithmetic, if the constraint is satisfied initially, it remains satisfied during the evolution. In a numerical evolution, however, truncation error and roundoff error generally produce small but nonzero values in $\mathcal{C}$. To that end we reformulate the evolution system so that constraint violations propagate in a controlled way.

We now introduce the auxiliary variable
\begin{equation}
\Gamma
\equiv
D_iA^i .
\end{equation}
This definition introduces an additional auxiliary constraint,
\begin{equation}
\mathcal{G}
\equiv
\Gamma-D_iA^i .
\end{equation}
For the original Maxwell system, $\mathcal{G}=0$. We can nevertheless promote $\Gamma$ to an independent evolved variable. Taking a time derivative of its defining relation gives
\begin{equation}
\partial_t\Gamma
=
D_i\partial_t A^i .
\end{equation}
Using the evolution equation for $A_i$, we find
\begin{equation}
\partial_t\Gamma
=
-D_iE^i
-
D_iD^i\Phi .
\end{equation}
Since $\mathcal{C}=0$ for constraint satisfying data, we may add an arbitrary multiple of $\mathcal{C}$ to the right hand side without changing the continuum solutions. Thus, we write
\begin{equation}
\partial_t\Gamma
=
-D_iE^i
-
D_iD^i\Phi
+
\eta\mathcal{C},
\end{equation}
where $\eta$ is a constant. Using
\begin{equation}
\mathcal{C}=D_iE^i-4\pi\rho ,
\end{equation}
this can also be written as
\begin{equation}
\partial_t\Gamma
=
(\eta-1)D_iE^i
-
D_iD^i\Phi
-
4\pi\eta\rho .
\end{equation}

With this new variable, the Maxwell evolution system becomes
\begin{align}
\partial_t A_i
&=
-E_i
-
D_i\Phi ,
\\
\partial_t E_i
&=
-D^2A_i
+
D_i\Gamma
-
4\pi J_i ,
\\
\partial_t\Gamma
&=
-D_iE^i
-
D_iD^i\Phi
+
\eta\mathcal{C}.
\end{align}
Here $D^2\equiv D_iD^i$. This system is equivalent to the original Maxwell system when the constraints are satisfied, but it changes how constraint violations behave.

To see this, start from
\begin{equation}
\mathcal{C}
=
D_iE^i
-
4\pi\rho .
\end{equation}
Taking a time derivative gives
\begin{equation}
\partial_t\mathcal{C}
=
D_i\partial_t E^i
-
4\pi\partial_t\rho .
\end{equation}
Using charge conservation,
\begin{equation}
\partial_t\rho
=
-D_i J^i ,
\end{equation}
we obtain
\begin{equation}
\partial_t\mathcal{C}
=
D_i\partial_t E^i
+
4\pi D_i J^i .
\end{equation}
Substituting the evolution equation for $E_i$ gives
\begin{align}
\partial_t\mathcal{C}
&=
D^i
\left(
-D^2A_i
+
D_i\Gamma
-
4\pi J_i
\right)
+
4\pi D_iJ^i
\\
&=
-D^iD^2A_i
+
D^iD_i\Gamma .
\end{align}
Assuming that the spatial derivatives commute, this becomes
\begin{equation}
\partial_t\mathcal{C}
=
-D^2D^iA_i
+
D^2\Gamma .
\end{equation}
Therefore,
\begin{equation}
\partial_t\mathcal{C}
=
D^2
\left(
\Gamma-D_iA^i
\right)
=
D^2\mathcal{G}.
\end{equation}

We next compute the evolution of $\mathcal{G}$. Taking a time derivative gives
\begin{equation}
\partial_t\mathcal{G}
=
\partial_t\Gamma
-
D_i\partial_t A^i .
\end{equation}
Using the evolution equations for $\Gamma$ and $A_i$, we find
\begin{align}
\partial_t\mathcal{G}
&=
\left(
-D_iE^i
-
D_iD^i\Phi
+
\eta\mathcal{C}
\right)
-
\left(
-D_iE^i
-
D_iD^i\Phi
\right)
\\
&=
\eta\mathcal{C}.
\end{align}
Combining this equation with
\begin{equation}
\partial_t\mathcal{C}
=
D^2\mathcal{G}
\end{equation}
gives
\begin{equation}
\partial_t^2\mathcal{C}
=
D^2\partial_t\mathcal{G}
=
\eta D^2\mathcal{C}.
\end{equation}
Equivalently,
\begin{equation}
\boxed{
\left(
-\partial_t^2+\eta D_iD^i
\right)
\mathcal{C}
=
0
} .
\end{equation}
Thus, for $\eta>0$, electric constraint violations propagate with speed $\sqrt{\eta}$. Analytically, this reformulation does not change the physical solutions of Maxwell's equations, because the added terms vanish for constraint satisfying data. Numerically, however, it changes the behavior of constraint violations by allowing them to propagate rather than remain fixed in place.

The BSSN formulation uses a related idea. It introduces auxiliary variables, most notably the conformal connection functions, and adds suitable multiples of the constraints to improve the hyperbolic structure and numerical behavior of the Einstein evolution system.

\subsubsection{BSSN Equations}

We begin with the ADM equations:
\begin{align}
R + K^2 - K_{ij}K^{ij}
&=
16\pi \rho,
&&\text{Hamiltonian constraint},
\\[0.5em]
D_j \left( K^{ij} - \gamma^{ij}K \right)
&=
8\pi j^i,
&&\text{Momentum constraint},
\\[0.5em]
\partial_t \gamma_{ij}
&=
-2\alpha K_{ij}
+
\mathcal{L}_{\beta}\gamma_{ij},
&&\text{Metric evolution},
\\[0.5em]
\partial_t K_{ij}
&=
\alpha
\left(
R_{ij}
- 2K_{ik}K^k{}_{j}
+ K K_{ij}
\right)
- D_i D_j \alpha
\notag \\
&\quad
- 4\pi \alpha
\left[
2S_{ij}
-
\gamma_{ij}(S-\rho)
\right]
+
\mathcal{L}_{\beta}K_{ij},
&&\text{Extrinsic curvature evolution}.
\end{align}

The Ricci tensor can be written in several equivalent forms. In terms of the spatial metric $\gamma_{ij}$, it contains second derivatives of the metric:
\begin{equation}
\begin{aligned}
R_{ij}
&=
-\frac{1}{2}\gamma^{kl}
\left(
\partial_k\partial_l\gamma_{ij}
+
\partial_i\partial_j\gamma_{kl}
-
\partial_i\partial_l\gamma_{kj}
-
\partial_k\partial_j\gamma_{il}
\right)
\\
&\quad
+
\gamma^{kl}
\left(
\Gamma^{m}{}_{il}\Gamma_{mkj}
-
\Gamma^{m}{}_{ij}\Gamma_{mkl}
\right).
\end{aligned}
\end{equation}
The term $\gamma^{kl}\partial_k\partial_l\gamma_{ij}$ acts like a Laplace operator on the metric. However, the remaining second derivative terms prevent the system from immediately taking a simple wave equation form. In analogy with the Maxwell example above, we introduce an auxiliary variable related to spatial derivatives of the metric.

For the physical spatial metric, define the contracted Christoffel symbol
\begin{equation}
\Gamma^i
\equiv
\gamma^{jk}\Gamma^i{}_{jk}.
\end{equation}
In terms of this quantity, the Ricci tensor may be written as
\begin{equation}
\begin{aligned}
R_{ij}
&=
-\frac{1}{2}\gamma^{kl}\partial_k\partial_l\gamma_{ij}
+
\gamma_{k(i}\partial_{j)}\Gamma^k
+
\Gamma^k\Gamma_{(ij)k}
\\
&\quad
+
2\gamma^{lm}\Gamma^k{}_{m(i}\Gamma_{j)kl}
+
\gamma^{kl}\Gamma^m{}_{il}\Gamma_{mkj}.
\end{aligned}
\end{equation}

The contracted Christoffel symbol is related to the divergence of the inverse metric by
\begin{equation}
g^{bc}\Gamma^a{}_{bc}
=
-\frac{1}{\sqrt{|g|}}
\partial_b
\left(
\sqrt{|g|}g^{ab}
\right).
\end{equation}
Applying this identity to the spatial metric gives
\begin{equation}
\label{eq:Gamma_eq}
\Gamma^i
=
-\frac{1}{\sqrt{\gamma}}
\partial_j
\left(
\sqrt{\gamma}\gamma^{ij}
\right).
\end{equation}

The BSSN formulation applies this idea to a conformally related metric. We introduce the conformal decomposition
\begin{equation}
\gamma_{ij}
=
e^{4\varphi}\tilde{\gamma}_{ij}.
\end{equation}
The conformal factor is chosen so that
\begin{equation}
\tilde{\gamma}=1,
\end{equation}
where $\tilde{\gamma}$ is the determinant of $\tilde{\gamma}_{ij}$. Taking the determinant of the conformal decomposition gives
\begin{equation}
\varphi
=
\frac{1}{12}\ln\gamma,
\end{equation}
where $\gamma$ is the determinant of the physical spatial metric $\gamma_{ij}$.

The conformal connection functions are then defined by
\begin{equation}
\tilde{\Gamma}^i
\equiv
\tilde{\gamma}^{jk}\tilde{\Gamma}^i{}_{jk}.
\end{equation}
Since $\tilde{\gamma}=1$, Eq.~\eqref{eq:Gamma_eq} gives the simpler expression
\begin{equation}
\tilde{\Gamma}^i
=
-\partial_j\tilde{\gamma}^{ij}.
\end{equation}
In the BSSN formulation, these $\tilde{\Gamma}^i$ are promoted to independent evolved variables. This is directly analogous to introducing $\Gamma=D_iA^i$ in the Maxwell example: the new variable agrees with a spatial derivative of the fundamental field when the auxiliary constraint is satisfied, but treating it independently allows the evolution system to be rewritten with improved hyperbolic properties.

The physical Ricci tensor is split into two pieces,
\begin{equation}
R_{ij}
=
\tilde{R}_{ij}
+
R^{\varphi}_{ij},
\end{equation}
where $\tilde{R}_{ij}$ is computed from the conformal metric $\tilde{\gamma}_{ij}$, and $R^{\varphi}_{ij}$ contains terms involving the conformal factor $\varphi$.

The conformal Ricci tensor is
\begin{equation}
\begin{aligned}
\tilde{R}_{ij}
&=
-\frac{1}{2}\tilde{\gamma}^{kl}
\partial_k\partial_l \tilde{\gamma}_{ij}
+
\tilde{\gamma}_{k(i}\partial_{j)}\tilde{\Gamma}^{k}
+
\tilde{\Gamma}^{k}\tilde{\Gamma}_{(ij)k}
\\
&\quad
+
2\tilde{\gamma}^{lm}
\tilde{\Gamma}^{k}{}_{m(i}\tilde{\Gamma}_{j)kl}
+
\tilde{\gamma}^{kl}
\tilde{\Gamma}^{m}{}_{il}\tilde{\Gamma}_{mkj}.
\end{aligned}
\end{equation}
Here $\tilde{\Gamma}^{i}{}_{jk}$ is the Christoffel symbol associated with the conformal metric $\tilde{\gamma}_{ij}$, and indices on conformal quantities are raised and lowered with $\tilde{\gamma}_{ij}$.

The conformal factor contribution is
\begin{equation}
\begin{aligned}
R^{\varphi}_{ij}
&=
-2
\left(
\tilde{D}_i\tilde{D}_j\varphi
+
\tilde{\gamma}_{ij}\tilde{D}^k\tilde{D}_k\varphi
\right)
\\
&\quad
+
4
\left[
\left(\tilde{D}_i\varphi\right)
\left(\tilde{D}_j\varphi\right)
-
\tilde{\gamma}_{ij}
\left(\tilde{D}^k\varphi\right)
\left(\tilde{D}_k\varphi\right)
\right],
\end{aligned}
\end{equation}
where $\tilde{D}_i$ is the covariant derivative compatible with $\tilde{\gamma}_{ij}$.

The contracted conformal Christoffel symbols are defined by
\begin{equation}
\tilde{\Gamma}^{i}
\equiv
\tilde{\gamma}^{jk}\tilde{\Gamma}^{i}{}_{jk}
=
-\frac{1}{\sqrt{\tilde{\gamma}}}
\partial_j
\left(
\sqrt{\tilde{\gamma}}\,
\tilde{\gamma}^{ij}
\right).
\end{equation}
Since the conformal metric is chosen to satisfy $\tilde{\gamma}=1$, this reduces to
\begin{equation}
\tilde{\Gamma}^{i}
=
-\partial_j\tilde{\gamma}^{ij}.
\end{equation}
In the BSSN formulation, the quantities $\tilde{\Gamma}^i$ are promoted to independent evolved variables.

The extrinsic curvature is decomposed into its trace and trace free parts:
\begin{equation}
K_{ij}
=
A_{ij}
+
\frac{1}{3}\gamma_{ij}K,
\end{equation}
where
\begin{equation}
K
=
\gamma^{ij}K_{ij}.
\end{equation}
The trace free part is then conformally rescaled according to
\begin{equation}
A_{ij}
=
e^{4\varphi}\tilde{A}_{ij},
\end{equation}
or equivalently,
\begin{equation}
\tilde{A}_{ij}
=
e^{-4\varphi}
\left(
K_{ij}
-
\frac{1}{3}\gamma_{ij}K
\right).
\end{equation}

The evolved BSSN variables are therefore
\begin{equation}
\varphi,\qquad
\tilde{\gamma}_{ij},\qquad
\tilde{A}_{ij},\qquad
K,\qquad
\tilde{\Gamma}^i .
\end{equation}
These variables are supplemented by the algebraic constraints
\begin{equation}
\tilde{\gamma}=1,
\qquad
\tilde{\gamma}^{ij}\tilde{A}_{ij}=0,
\end{equation}
and by the auxiliary connection constraint
\begin{equation}
\mathcal{G}^i
\equiv
\tilde{\Gamma}^i
+
\partial_j\tilde{\gamma}^{ij}
=
0 .
\end{equation}
Together with appropriate gauge conditions for the lapse and shift, these variables form the standard BSSN formulation.

\noindent\textbf{The BSSN equations are}
\begin{equation}
\boxed{
\begin{aligned}
\partial_t \varphi
&=
-\frac{1}{6}\alpha K
+
\beta^i \partial_i \varphi
+
\frac{1}{6}\partial_i \beta^i ,
&&\text{Conformal exponent evolution},
\\[0.75em]
%
\partial_t K
&=
-\gamma^{ij}D_iD_j\alpha
+
\alpha
\left(
\tilde{A}_{ij}\tilde{A}^{ij}
+
\frac{1}{3}K^2
\right)
+
4\pi\alpha(\rho+S)
+
\beta^i\partial_i K ,
&&\text{Trace evolution},
\\[0.75em]
%
\partial_t \tilde{\gamma}_{ij}
&=
-2\alpha\tilde{A}_{ij}
+
\beta^k\partial_k\tilde{\gamma}_{ij}
+
\tilde{\gamma}_{ik}\partial_j\beta^k
+
\tilde{\gamma}_{kj}\partial_i\beta^k
-
\frac{2}{3}\tilde{\gamma}_{ij}\partial_k\beta^k ,
&&\text{Conformal metric evolution},
\\[0.75em]
%
\partial_t \tilde{A}_{ij}
&=
e^{-4\varphi}
\left[
-D_iD_j\alpha
+
\alpha
\left(
R_{ij}
-
8\pi S_{ij}
\right)
\right]^{\mathrm{TF}}
\\
&\quad
+
\alpha
\left(
K\tilde{A}_{ij}
-
2\tilde{A}_{il}\tilde{A}^{l}{}_{j}
\right)
\\
&\quad
+
\beta^k\partial_k\tilde{A}_{ij}
+
\tilde{A}_{ik}\partial_j\beta^k
+
\tilde{A}_{kj}\partial_i\beta^k
-
\frac{2}{3}\tilde{A}_{ij}\partial_k\beta^k ,
&&\text{Trace free curvature evolution},
\\[0.75em]
%
\partial_t \tilde{\Gamma}^{i}
&=
-2\tilde{A}^{ij}\partial_j\alpha
+
2\alpha
\left(
\tilde{\Gamma}^{i}{}_{jk}\tilde{A}^{kj}
-
\frac{2}{3}\tilde{\gamma}^{ij}\partial_j K
-
8\pi\tilde{\gamma}^{ij}j_j
+
6\tilde{A}^{ij}\partial_j\varphi
\right)
\\
&\quad
+
\beta^j\partial_j\tilde{\Gamma}^{i}
-
\tilde{\Gamma}^{j}\partial_j\beta^i
+
\frac{2}{3}\tilde{\Gamma}^{i}\partial_j\beta^j
\\
&\quad
+
\frac{1}{3}\tilde{\gamma}^{li}\partial_l\partial_j\beta^j
+
\tilde{\gamma}^{lj}\partial_j\partial_l\beta^i ,
&&\text{Conformal connection evolution}.
\end{aligned}
}
\end{equation}

\subsubsection{Gauge and Slicing Conditions}

The BSSN equations do not fully determine the evolution until the coordinate gauge is specified. In particular, one must choose evolution conditions for the lapse $\alpha$ and the shift $\beta^i$. These choices determine how the spatial slices are placed in spacetime and how the spatial coordinates move from one slice to the next. We summarize several common choices below.

\paragraph{Geodesic slicing.}
A simple choice is
\begin{equation}
\alpha=1,
\qquad
\beta^i=0 .
\end{equation}
This is called geodesic slicing. Although this gauge is simple, it is usually a poor choice for numerical relativity simulations of black holes. The acceleration of the normal observers is
\begin{equation}
a_a
=
n^b\nabla_b n_a
=
D_a\ln\alpha .
\end{equation}
For geodesic slicing, $\alpha=1$, so $a_a=0$. Thus, the normal observers follow geodesics. In a strong gravitational field, these freely falling observers can focus and form coordinate singularities, making geodesic slicing unsuitable for most black hole evolutions.

\paragraph{Maximal slicing.}
Another possibility is to control the convergence of the normal observers by controlling the mean curvature of the spatial slices. In maximal slicing, one imposes
\begin{equation}
K=0 .
\end{equation}
If this condition is maintained throughout the evolution, then the slicing condition requires
\begin{equation}
\partial_t K = 0,
\end{equation}
up to shift terms. Using the trace evolution equation, this condition gives an elliptic equation for the lapse. Schematically,
\begin{equation}
D^iD_i\alpha
=
\alpha
\left(
K_{ij}K^{ij}
+
4\pi(\rho+S)
\right),
\end{equation}
for $K=0$. Maximal slicing has several desirable properties, including singularity avoidance. In particular, the lapse tends to collapse toward zero near strong field regions, causing those parts of the slice to advance more slowly in coordinate time. This behavior is known as collapse of the lapse. However, maximal slicing also has disadvantages, such as grid stretching, which can produce sharp coordinate gradients that are difficult to resolve numerically.

\paragraph{Harmonic coordinates.}
Harmonic coordinates are defined by requiring the spacetime coordinates to satisfy wave equations,
\begin{equation}
\square x^a = 0 .
\end{equation}
Equivalently, one may impose
\begin{equation}
\Gamma^a
=
g^{bc}\Gamma^a{}_{bc}
=
0 .
\end{equation}
A related condition is harmonic slicing, in which only the time component of this condition is imposed. If we also set $\beta^i=0$, then
\begin{equation}
\Gamma^0
=
-\frac{1}{\alpha\sqrt{\gamma}}
\partial_t
\left(
\frac{\sqrt{\gamma}}{\alpha}
\right)
=
0 .
\end{equation}
This gives the lapse evolution equation
\begin{equation}
\partial_t \alpha
=
-\alpha^2 K .
\end{equation}
Harmonic slicing is closely related to wave coordinates and often leads to well behaved hyperbolic evolution systems. However, modified harmonic slicing conditions are typically more useful for black hole evolutions.


\paragraph{\boldmath $1+\log$ slicing and the Gamma driver condition.}
A common generalization of harmonic slicing is the Bona--Masso family of slicing conditions,
\begin{equation}
\partial_t \alpha
=
-\alpha^2 f(\alpha)K .
\end{equation}
For $f(\alpha)=1$, this reduces to harmonic slicing. For $f(\alpha)=0$, the lapse is constant in time, giving geodesic slicing when $\alpha=1$. Larger values of $f(\alpha)$ make the lapse respond more strongly to the local value of $K$.

A particularly useful choice is
\begin{equation}
f(\alpha)
=
\frac{2}{\alpha}.
\end{equation}
This gives
\begin{equation}
\partial_t\alpha
=
-2\alpha K,
\end{equation}
which is the standard $1+\log$ slicing condition. This choice has singularity avoidance properties similar to maximal slicing while retaining good behavior in nearly flat regions.

When the shift is nonzero, it is common to include an advective term and write
\begin{equation}
\left(\partial_t-\beta^i\partial_i\right)\alpha
=
-\alpha^2 f(\alpha)K .
\end{equation}
For $f(\alpha)=2/\alpha$, this becomes
\begin{equation}
\left(\partial_t-\beta^i\partial_i\right)\alpha
=
-2\alpha K .
\end{equation}

We also need an evolution condition for the shift. A common choice is the Gamma driver condition,
\begin{equation}
\left(\partial_t-\beta^j\partial_j\right)\beta^i
=
\mu B^i,
\end{equation}
together with
\begin{equation}
\left(\partial_t-\beta^j\partial_j\right)B^i
=
\left(\partial_t-\beta^j\partial_j\right)\tilde{\Gamma}^i
-
\eta B^i .
\end{equation}
Here $B^i$ is an auxiliary variable, while $\mu$ and $\eta$ are parameters that control the characteristic speed and damping of the shift condition. The combination of $1+\log$ slicing and the Gamma driver shift condition is called the moving puncture gauge. A major advantage of this gauge is that it allows black hole spacetimes to be evolved without excising the singularity from the computational domain.

The parameter $\mu$ controls how quickly the shift responds to changes in the conformal connection functions, while $\eta$ damps the auxiliary variable $B^i$ and suppresses oscillatory behavior in the shift.
